% GENERATED by paper/tools/exp/countmap-figs.mjs -- do not edit by hand.
\begin{tikzpicture}[x=1cm, y=1cm, font=\scriptsize, >={Stealth[length=3pt]}]
  % Row 1: one family, its count, its fractional part.
  \node[anchor=south west, inner sep=0] at (0.000,0) {
\includegraphics[width=4.2cm]{countmap-family.png}};
  \draw[soft] (0.000,0) rectangle (4.200,4.2);
  \node[anchor=south west, inner sep=0] at (4.650,0) {\includegraphics[width=4.2cm]{countmap-count.png}};
  \draw[soft] (4.650,0) rectangle (8.850,4.2);
  \node[anchor=south west, inner sep=0] at (9.300,0) {\includegraphics[width=4.2cm]{countmap-frac.png}};
  \draw[soft] (9.300,0) rectangle (13.500,4.2);
  \node[accent, fill=white, inner sep=0.6pt] at (0.173,4.000) {0};
  \node[accent, fill=white, inner sep=0.6pt] at (0.635,4.000) {1};
  \node[accent, fill=white, inner sep=0.6pt] at (1.097,4.000) {2};
  \node[accent, fill=white, inner sep=0.6pt] at (1.559,4.000) {3};
  \node[accent, fill=white, inner sep=0.6pt] at (2.021,4.000) {4};
  \node[accent, fill=white, inner sep=0.6pt] at (2.483,4.000) {5};
  \node[accent, fill=white, inner sep=0.6pt] at (2.945,4.000) {6};
  \node[accent, fill=white, inner sep=0.6pt] at (3.407,4.000) {7};
  \node[accent, fill=white, inner sep=0.6pt] at (3.869,4.000) {8};
  \node[anchor=north, align=center, text width=4.2cm] at (2.100,-0.080) {(a) a family, members numbered};
  \node[anchor=north, align=center, text width=4.2cm] at (6.750,-0.080) {(b) its count $\xi(x)$, a function on the plane; the members are its whole-number level sets};
  \node[anchor=north, align=center, text width=4.2cm] at (11.400,-0.080) {(c) where you are in the member: $\xi\bmod1$, a ramp that jumps at each member};
  % Row 2: two families read at one point, the state on the torus, the drawing recomputed from the state.
  \begin{scope}[shift={(0,-5.9)}]
  \node[anchor=south west, inner sep=0] at (0.000,0) {\includegraphics[width=4.2cm]{countmap-pair.png}};
  \draw[soft] (0.000,0) rectangle (4.200,4.2);
    \draw[accent, ->, line width=0.7pt] (2.730,2.339) -- (3.192,2.029);
    \fill[accent] (2.961,2.184) circle (1.6pt);
    \node[accent, anchor=south west, fill=white, inner sep=1pt] at (3.041,2.264) {$x$};
  \node[anchor=north, align=center, text width=4.2cm] at (2.100,-0.300) {(d) two families, drawn. At $x$: $\xi_1=1.72$, $\xi_2=3.80$; a short step along the arrow advances both};
    % The torus: the unit square with its edges glued, the picture, the state, the path.
    \begin{scope}[shift={(4.650,0)}]
      % The picture I: family 1's stroke is a vertical band about u1 = 0 (blue), family 2's a horizontal band about u2 = 0 (orange).
      \fill[cool!45] (0,0) rectangle (0.286,4.2);
      \fill[cool!45] (3.914,0) rectangle (4.2,4.2);
      \fill[warm!45] (0,0) rectangle (4.2,0.315);
      \fill[warm!45] (0,3.885) rectangle (4.2,4.2);
      \foreach \ax/\bx in {0/0.286, 3.914/4.2} { \foreach \ay/\by in {0/0.315, 3.885/4.2} { \fill[ink!45] (\ax,\ay) rectangle (\bx,\by); } }
      % Level lines of the slow recipe (1,-1), faint, clipped to the square.
      \begin{scope}
        \clip (0,0) rectangle (4.2,4.2);
        \foreach \c in {-0.8,-0.6,...,0.81} { \draw[accent!35, thin] (\c*4.2,0) -- (\c*4.2+4.2,4.2); }
      \end{scope}
      \draw[ink] (0,0) rectangle (4.2,4.2);
      % Gluing arrows: single on left and right, double on top and bottom.
      \draw[ink, ->] (0,1.890) -- (0,2.310);
      \draw[ink, ->] (4.2,1.890) -- (4.2,2.310);
      \draw[ink, ->] (1.806,0) -- (2.226,0);
      \draw[ink, ->] (1.806,4.2) -- (2.226,4.2);
      \draw[ink, ->] (1.974,0) -- (2.394,0);
      \draw[ink, ->] (1.974,4.2) -- (2.394,4.2);
      % The path of the segment through x.
      \draw[accent, line width=1.1pt] (0.664,1.443) -- (0.700,1.483) -- (0.737,1.524) -- (0.775,1.564) -- (0.813,1.605) -- (0.852,1.646) -- (0.892,1.686) -- (0.933,1.728) -- (0.974,1.769) -- (1.017,1.810) -- (1.060,1.852) -- (1.104,1.894) -- (1.149,1.935) -- (1.195,1.977) -- (1.241,2.020) -- (1.289,2.062) -- (1.337,2.104) -- (1.386,2.147) -- (1.436,2.190) -- (1.487,2.233) -- (1.538,2.276) -- (1.591,2.319) -- (1.644,2.363) -- (1.698,2.406) -- (1.753,2.450) -- (1.809,2.493) -- (1.865,2.537) -- (1.922,2.581) -- (1.980,2.626) -- (2.039,2.670) -- (2.099,2.714) -- (2.159,2.759) -- (2.220,2.804) -- (2.282,2.849) -- (2.345,2.894) -- (2.408,2.939) -- (2.472,2.984) -- (2.537,3.029) -- (2.602,3.075) -- (2.668,3.120) -- (2.735,3.166) -- (2.802,3.212) -- (2.870,3.258) -- (2.938,3.304) -- (3.007,3.350) -- (3.077,3.397) -- (3.147,3.443) -- (3.218,3.490) -- (3.289,3.536) -- (3.360,3.583) -- (3.432,3.630) -- (3.505,3.677) -- (3.577,3.724) -- (3.650,3.771) -- (3.724,3.819) -- (3.798,3.866) -- (3.872,3.914) -- (3.947,3.961) -- (4.021,4.009) -- (4.096,4.057) -- (4.171,4.105);
      \draw[accent, line width=1.1pt] (0.047,4.153);
      \draw[accent, line width=1.1pt] (0.122,0.001) -- (0.198,0.049) -- (0.274,0.098) -- (0.350,0.146) -- (0.426,0.195) -- (0.502,0.243) -- (0.578,0.292) -- (0.654,0.341) -- (0.730,0.390) -- (0.806,0.439) -- (0.882,0.488) -- (0.958,0.537) -- (1.033,0.587) -- (1.109,0.636) -- (1.184,0.685) -- (1.259,0.735) -- (1.334,0.785) -- (1.409,0.834) -- (1.483,0.884) -- (1.557,0.934) -- (1.631,0.984) -- (1.704,1.034) -- (1.778,1.084) -- (1.850,1.135) -- (1.922,1.185) -- (1.994,1.235) -- (2.066,1.286);
      \fill[accent] (3.007,3.350) circle (1.8pt);
      \node[accent, anchor=south east, fill=white, inner sep=1pt] at (2.947,3.410) {$\Phi(x)=(0.72,0.80)$};
      \node[anchor=north, font=\tiny] at (2.1,-0.03) {$\xi_1\bmod1$};
      \node[anchor=south, rotate=90, font=\tiny] at (-0.03,2.1) {$\xi_2\bmod1$};
    \end{scope}
  \node[anchor=north, align=center, text width=4.2cm] at (6.750,-0.300) {(e) the state $\Phi(x)=(\xi_1,\xi_2)\bmod1$ on the torus: the square with its edges glued (arrows). Blue: states where the lines have ink; orange: the rings. Red: the image of the arrow in (d). Faint: $\xi_1-\xi_2$ constant};
  \node[anchor=south west, inner sep=0] at (9.300,0) {\includegraphics[width=4.2cm]{countmap-composite.png}};
  \draw[soft] (9.300,0) rectangle (13.500,4.2);
  \node[anchor=north, align=center, text width=4.2cm] at (11.400,-0.300) {(f) the drawing computed from the state alone: at every pixel look up $\Phi$, then read the picture in (e). It is (d) again: $S=I\circ\Phi$};
  \end{scope}
\end{tikzpicture}
